\section{Reduction of the deeper step to a local Poisson moment}\label{sec:moment}

We now derive the conditional reduction of BB Note15~\cite{BB15},
with every part of the boundary included.

\subsection{A collar with a fixed Green margin}

Fix $X$ large and hold $t=t_X$ throughout this section. Set
\[
 \ell=17/32,\quad u=39/64,\quad b=37/64,\quad a=73/128,
 \quad d_s=u-\ell=5/64,\quad W=\frac{3d_s}{\pi}\log X,
\]
\[
 D_X=\{|\Re z-X|<W,\ \ell<\Im z<u\},\quad
 p=X+\ii b,\quad K_X=\{X+\ii y:9/16\le y\le a\}.
\]
The already acquired centerline enclosure is below $a$.
The disk centered at $X+\ii a$ of radius $5/128$ lies in $D_X$
eventually. In units $1/128$, the observation and deepest target have
coordinates $+1,-1$ in a disk of radius five. Formula
\eqref{eq:disk-cost} and domain monotonicity give
\begin{equation}\label{eq:green-signal}
 \inf_{z\in K_X}G_{D_X}(p,z)\ge\log(13/5).
\end{equation}

Normalize by the same independently nonzero carrier,
\[
 F(z)=\frac{H_{t_X}(z)}{C_-(t_X,z)\zeta(s_{t_X}(z))}.
\]
At every point of the closed rectangle recenter the fixed provider
chart at its real coordinate. The time remains $t_X$.
Since $|t_XL_x-1|=O(W/(XL_X))$, an eventual gate gives
\begin{equation}\label{eq:carrier-floor-66}
 \Re s_{t_X}\ge65/64-1/40000>1+1/65,
 \qquad |1/\zeta(s_{t_X})|\le66.
\end{equation}
The denominator is analytic and nonzero on a neighborhood of the
rectangle. At the observation and on the top, the complete comparison
of Section~\ref{sec:width} applies with the same small clock drift:
\begin{equation}\label{eq:anchor}
 |F(p)|\ge15/16,\qquad F=1+o(1)\quad\hbox{uniformly on the top}.
\end{equation}
Only the top and observation use that estimate; it is not assumed
on the deeper bottom edge.

The quadratic consumer excludes $K_X$ if its complete boundary
payment is at most $9/2$, because
\begin{equation}\label{eq:margin-81}
 (15/16)^2\{(13/5)^2-1\}=81/16>9/2.
\end{equation}
Taking $\kappa_s=\pi/(2d_s)$ in the end barrier shows that both
sides together have mass $O(X^{-3/2})$. The complete provider and
positive finite population give $|H_{t_X}/C_-|\ll X^{1/4}$ there.
After~\eqref{eq:carrier-floor-66}, the side quadratic cost is
$O(X^{-1})$. The top costs $o(1)$ by~\eqref{eq:anchor}.
It remains to pay the bottom at its actual harmonic measure.

\subsection{The kernel and the complete bottom-edge reduction}

The lower-edge Poisson kernel of the containing strip at $p$ is
\begin{equation}\label{eq:poisson}
 P(v)=\frac{\sin(3\pi/5)}
 {2d_s\{\cosh(\pi v/d_s)-\cos(3\pi/5)\}},\qquad d_s=5/64.
\end{equation}
It is positive and even. With the Fourier convention
$\widehat P(\omega)=\int P(v)\ee^{\ii\omega v}\dd v$, one has
\begin{equation}\label{eq:poisson-transform}
 \widehat P(\omega)=\frac{\sinh(|\omega|/32)}
                         {\sinh(5|\omega|/64)},\qquad
 \widehat P(0)=\int P(v)\dd v=2/5.
\end{equation}
To derive the transform, extend the lower-boundary datum
$\ee^{\ii\omega v}$ harmonically into the strip with zero datum on
the top. Separation of variables gives height factor
$\sinh((u-b)|\omega|)/\sinh(d_s|\omega|)$.
At zero frequency the extension is linear in height. Equivalently,
the substitution $s=\tanh(\pi v/(2d_s))$ evaluates the integral
of~\eqref{eq:poisson} directly. Also $P(v)\ll\ee^{-\pi|v|/d_s}$
for $|v|\ge1$.

For any nonnegative bottom datum on the rectangle, extend it by zero
to the rest of the strip bottom. Its strip extension is nonnegative
on both rectangle sides and dominates the rectangle extension by the
maximum principle. Thus the rectangle's bottom harmonic measure is
dominated by $P(v)\dd v$ on $|v|\le W$.

We now justify the exact tail in~\eqref{eq:raw-tail-intro}.
Split the dominant arm at $K=\lfloor\ee^{L_X/32}\rfloor$.
On the finite head the weighted summability exponent exceeds
$65/64-1/40000-1/128>1+1/256$, so heat removal costs $O(1/L_X)$.
The ordinary zeta tail costs $O(X^{-1/2080})$ using
\eqref{eq:carrier-floor-66}.
At $y=17/32$, the reflected primary power with $(1/6,2/3)$
has endpoints $-19/192$ and $-11/384$ at $c=1$.
The clock drift can be made to cost less than $1/160000$;
the complete reflected sum is consequently
$O((1+L_X)X^{-1/48})$, including its inverse-frequency terms.
The correction and full remainder in~\eqref{eq:ratio-correction}
also tend to zero uniformly.

At $x=X+v$, $|v|\le W$, the natural upper cutoff
$\lfloor\sqrt{x/(4\pi)}\rfloor$ differs from the one at $X$
by at most one eventually. The boundary term costs $O(X^{-57/128})$.
For common terms the ratio to the frozen summand in $T_X(v)$ is
\[
 \exp\left(-\frac{t_X}{4}(L(X+v+\ii\ell)-L_X)\log n\right)
 =1+O(W/X).
\]
The complete absolute population is $O(X^{1/4})$, so the summed
freezing error is $o(1)$. We have proved, uniformly along the bottom,
\begin{equation}\label{eq:bottom-tail}
 F(X+v+\ii\ell)-1
 =\frac{T_X(v)+o(1)}{\zeta(s_{t_X}(X+v+\ii\ell))}.
\end{equation}
All heat-weighted sums are finite, all cutoffs are actual integer
cutoffs, and the entire correction and source remainder have been
included.

If $Q(X)\le4/66^2$, Poisson domination and the triangle inequality
in $L^2(P(v)\dd v)$ give bottom cost at most
\[
 66^2(\sqrt{Q(X)}+o(1))^2\le4+o(1).
\]
Adding the top and both sides gives at most $9/2$ eventually.
Equations~\eqref{eq:green-signal}--\eqref{eq:margin-81} exclude $K_X$.
The acquired centerline enclosure excludes the portion above $a$,
and conjugation covers the lower half-plane. This proves
Proposition~\ref{prop:intro-green}, conditional on its stated
every-center premise.

\subsection{Why signs and exact coefficients remain essential}

Expanding the finite squared norm yields the exact signed form
\begin{equation}\label{eq:Q-pairs}
 Q(X)=\frac25\sum_{K<n\le N}w_n^2
 +2\Re\sum_{K<m<n\le N}w_mw_n\ee^{\ii\tau_X\log(n/m)}
              \widehat P\left(\tfrac12\log(n/m)\right).
\end{equation}
The diagonal is $O(K^{-25/32})=O(X^{-25/1024})$, since
$w_n^2\le n^{-57/32}$. The issue is the signed off-diagonal sum.
For $N/2<m<n\le N$, the weights are comparable to $X^{-57/128}$
and the multiplier is bounded below by a positive constant.
The termwise absolute allowance from those pairs alone is therefore
of order at least $X^{7/64}$. This says nothing about a lower bound
for the signed canonical value.

The same issue appears in a coefficient-independent Gram estimate.
On a dyad $(D,2D]$, the unphased matrix with entries
$\widehat P(\log(n/m)/2)$ has all entries bounded below by a
positive constant. Its all-ones Rayleigh quotient is $\gg D$.
The center phase is a diagonal unitary conjugation, preserving its
operator norm. Multiplying that norm by $\sum w_n^2$ therefore
incurs a scale at least $D^2w_D^2$. If $D=\ee^{aL_X}$, its power
is $a^2/2-a/32$, positive when $a>1/16$.
This is a limitation of the coefficient-independent norm bound.

There is also an exact Gaussian representation of each finite arm.
For $s=65/64-\ii(\tau_X+v/2)$ and
$D_{A,B}(s)=\sum_{A<n\le B}n^{-s}$,
\begin{equation}\label{eq:weight-gaussian}
 \sum_{A<n\le B}w_n(X)\ee^{\ii(\tau_X+v/2)\log n}
 =\sqrt{L_X/\pi}\int_\R\ee^{-L_Xr^2}D_{A,B}(s-r)\dd r.
\end{equation}
This is the real Gaussian moment identity applied to a finite sum.
On a dyad of scale $\ee^{aL_X}$, absolute integration has exponent
$(r-1/64)a-r^2$. Its saddle $r=a/2$ gives $a^2/4-a/64$,
again positive for $a>1/16$. Applying a separate exponent-pair bound
under the integral returns the same weighted block exponent after
square completion. An infinite ordinary series is available only
where $\Re(s-r)>1$. Elsewhere a continuation argument must retain
the exact finite-cutoff remainder; the identity itself supplies no
new cancellation. These scoped comparisons explain why a signed
estimate of the actual moment is still needed.
